\documentclass[11pt]{article}
\usepackage[letterpaper,margin=1in]{geometry}
\usepackage[T1]{fontenc}
\usepackage[utf8]{inputenc}
\usepackage{lmodern,microtype}
\usepackage{amsmath,amssymb,mathtools,bm,slashed}
\usepackage{xcolor,enumitem,booktabs,array,needspace}
\usepackage[most]{tcolorbox}
\usepackage{fancyhdr,titlesec,parskip,hyperref}
\definecolor{qftblue}{HTML}{245A8D}
\definecolor{qftgreen}{HTML}{2F6B4F}
\definecolor{qftgray}{HTML}{F4F6F8}
\hypersetup{colorlinks=true,linkcolor=qftblue,urlcolor=qftblue,pdfauthor={Nima Lashkari}}
\pagestyle{fancy}\fancyhf{}
\lhead{PHYS 662: Quantum Field Theory I}\cfoot{\thepage}
\setlength{\headheight}{14pt}
\titleformat{\section}{\Large\bfseries\color{qftblue}}{\thesection.}{0.6em}{}
\titleformat{\subsection}{\large\bfseries\color{qftgreen}}{\thesubsection}{0.6em}{}
\setlist[enumerate]{leftmargin=2em,itemsep=0.25em,topsep=0.35em}
\setlist[itemize]{leftmargin=2em,itemsep=0.2em,topsep=0.35em}
\newcommand{\dd}{\mathrm d}
\newcommand{\ii}{\mathrm i}
\newcommand{\eps}{\varepsilon}
\newcommand{\tr}{\operatorname{tr}}
\newcommand{\Tr}{\operatorname{Tr}}
\newcommand{\diag}{\operatorname{diag}}
\newcommand{\Lag}{\mathcal L}
\newcommand{\D}{\mathcal D}
\newcommand{\vac}{\lvert0\rangle}
\newcommand{\vbra}{\langle0\rvert}
\newcommand{\Tord}{\mathrm T}
\newcommand{\ket}[1]{\lvert#1\rangle}
\newcommand{\bra}[1]{\langle#1\rvert}
\newtcolorbox{keyidea}[1][]{enhanced,breakable,colback=qftgray,colframe=qftblue,
 boxrule=0.7pt,arc=1.5mm,title={#1},fonttitle=\bfseries,coltitle=white}
\newtcolorbox{remarkbox}[1][]{enhanced,breakable,colback=white,colframe=qftgreen,
 boxrule=0.6pt,arc=1.5mm,title={#1},fonttitle=\bfseries,coltitle=white}
\newcommand{\intp}{\int\!\frac{\dd^4p}{(2\pi)^4}}
\newcommand{\MSbar}{\overline{\mathrm{MS}}}
\newcommand{\module}[1]{\clearpage\section{#1}}
\rhead{Homework Set 7}
\hypersetup{pdftitle={PHYS 662 Homework Set 7: Interactions and One Loop}}
\newlist{parts}{enumerate}{1}
\setlist[parts]{label=\textbf{(\alph*)},leftmargin=2.3em,itemsep=0.65em,topsep=0.5em}
\newcommand{\problem}[2]{\clearpage\section*{\color{qftblue}Problem #1: #2}}
\newcommand{\connection}[1]{{\small\textit{Lecture connection: #1}\par}\medskip}
\begin{document}
\begin{center}
{\scshape\Large PHYS 662 --- Quantum Field Theory I\par}
\vspace{0.3em}{\Huge\bfseries\color{qftblue}Homework Set 7\par}
\vspace{0.3em}{\large Interactions and Scalar Theory at One Loop\par}
\vspace{0.3em}{\large Lectures 21--22 \quad Fall 2026\par}
\end{center}\medskip
\textbf{Name:} \rule{0.43\textwidth}{0.4pt}\hfill
\textbf{Due:} \rule{0.24\textwidth}{0.4pt}

\textbf{Instructions.} Complete all four problems, worth 25 points each.
Show the derivation of each requested result, distinguish connected
from 1PI graphs, and display symmetry factors. Label external legs
and loop momenta on diagrams. Retain the Feynman prescription through
analytic continuation. Results from earlier parts may be reused.

\begin{keyidea}[Common conventions]
Use \(\hbar=c=1\), \(\eta=\diag(1,-1,-1,-1)\), the symmetric massive
vacuum, and
\[
 \Lag=\tfrac12(\partial\phi)^2-\tfrac12m^2\phi^2-\tfrac{\lambda}{4!}\phi^4,
 \quad m^2>0,\quad\lambda>0,\qquad
 D_F(p)=\frac{\ii}{p^2-m^2+\ii0}.
\]
The line already includes \(\ii\). Set \(Z[0]=1\),
\(W[J]=-\ii\log Z[J]\), and
\(G_n=\left.\ii^{-n}\delta_J^n Z\right|_0\).
Use \(Z_0[J]=\exp(-\tfrac12 JD_FJ)\), with spacetime integrations
implicit. Regulate coincident contractions before manipulating them.

For loop calculations use \(d=4-2\eps\) and
\(1/\bar\eps=1/\eps-\gamma_E+\log4\pi\). The counterterms are
\[
 \Lag_{\rm ct}=\frac{\delta Z_\phi}{2}(\partial\phi)^2
 -\frac{\delta m^2}{2}\phi^2
 -\frac{\mu^{2\eps}\delta\lambda}{4!}\phi^4.
\]
Write a two-point insertion as \(-\ii\Sigma_R\), where
\(\Sigma_R=\Sigma_{\rm loop}+\delta m^2-\delta Z_\phi p^2\).
Use \(M\) for the physical pole mass and \(m(\mu)\) for the
\(\MSbar\) mass. Remove the overall \(\mu^{2\eps}\) before quoting
the four-dimensional four-point amplitude.

Normalize particle states by
\(\langle\bm p|\bm q\rangle=2E_{\bm p}(2\pi)^3\delta^{(3)}(\bm p-\bm q)\),
and define \(\langle f|S-1|i\rangle=\ii(2\pi)^4
\delta^{(4)}(P_f-P_i)\mathcal M_{fi}\).
\end{keyidea}

\textbf{Coverage}
\begin{enumerate}
\item Source perturbation theory, normalization, and connected functions.
\item Diagram counting, momentum-space rules, and amputation.
\item The tadpole, regulators, and mass renormalization.
\item The bubble, coupling renormalization, running, and unitarity.
\end{enumerate}
Each problem is cumulative. The six-point tree in Problem 2 extends
the lecture's four-point examples without introducing a new interaction.

\problem{1}{From the Gaussian to interacting correlators}
\connection{Lecture 21, Sections 1--5.}
Let \(f(z)=\int\dd^4x\,D_F(z-x)J(x)\). In this problem do not
evaluate the divergent constant \(D_F(0)\); keep the same regulator
in numerator and denominator.
\begin{parts}
\item \textbf{(4 points)} Starting from the normalized interacting
path integral, derive
\[
 Z[J]=\frac{\exp\!\left[-\dfrac{\ii\lambda}{4!}\int\dd^4z
 (\delta/\ii\delta J(z))^4\right]Z_0[J]}
 {\left.\exp\!\left[-\dfrac{\ii\lambda}{4!}\int\dd^4z
 (\delta/\ii\delta J(z))^4\right]Z_0[J]\right|_{J=0}}.
\]
Explain why the denominator has no source dependence and why it must
be expanded to the same order as the numerator.

\item \textbf{(5 points)} Compute
\(Z_0^{-1}(\delta/\ii\delta J(z))^4Z_0\) explicitly. Identify the
terms with zero, two, and four sources and explain their combinatorial
coefficients. Expand the ratio through \(O(\lambda)\) and demonstrate
the cancellation of the source-independent vacuum term.

\item \textbf{(5 points)} Extract \(G_2(x,y)\) through
\(O(\lambda)\). Check it by applying \(\Box_x+m^2\) and using
\[
 (\Box_x+m^2)G_2(x,y)+\frac{\lambda}{3!}G_4(x,x,x,y)
 =-\ii\delta^{(4)}(x-y)
\]
with the free value of \(G_4\) in the second term. Track every factor
of \(\ii\), including the contact term.

\item \textbf{(4 points)} Extract the connected four-point function
at its leading nonzero order. Write the full four-point function
through \(O(\lambda)\) as its connected part plus three products of
full two-point functions, expanded consistently. Explain why
normalizing \(Z\) does not remove these products.

\item \textbf{(4 points)} Count contractions directly for the
one-vertex vacuum graph, the two-point tadpole, and the four-point
tree. Divide by the interaction factorial and give the three symmetry
factors. Keep the external labels fixed throughout the count.

\item \textbf{(3 points)} Show that
\(\left.\delta_J^nW\right|_0=\ii^{n-1}G_{n,c}\).
Use \(n=4\) to explain the subtraction in part (d). State why the
exact one-point Schwinger--Dyson equation does not generally imply
\(V'(\langle\phi\rangle)=0\), where
\(V(\phi)=m^2\phi^2/2+\lambda\phi^4/4!\).
\end{parts}

\problem{2}{Counting diagrams and amputating external legs}
\connection{Lecture 21, Sections 4--8.}
Use distinct labels for all external momenta and take them all incoming.
A reduced momentum-space correlator has its overall
\((2\pi)^4\delta^{(4)}(\sum_a p_a)\) removed. In parts (b)--(c), the
external momenta may be off shell and should be generic.
\begin{parts}
\item \textbf{(4 points)} For a connected quartic graph with \(E\)
external legs, \(I\) internal lines, \(V\) vertices, and \(L\)
independent loops, derive
\(4V=2I+E\) and \(L=I-V+1\). Determine \(V,I,L\) for the leading
connected six-point function and for the one-loop four-point function.
Explain why one loop does not mean one power of \(\lambda\).

\item \textbf{(4 points)} Enumerate all distinct six-point tree graphs
by partitioning the six external labels between the two vertices.
Explain why exchanging the two groups does not give a new graph.
Find the weight of one labeled graph by counting Wick contractions
in the order-\(\lambda^2\) source expansion, including its \(1/2!\).

\item \textbf{(5 points)} Write the amputated connected six-point
kernel as a sum over those partitions. For a three-element subset
\(A\), define \(P_A=\sum_{a\in A}p_a\). Display the internal
propagator and the overall phase in each term. Restore the external
propagators to obtain the reduced correlator. Are these trees 1PI?
What further conditions are needed to interpret the amputated result
as a physical \(3\to3\) amplitude?

\item \textbf{(5 points)} For the four-point partition \((12|34)\),
count the contractions that join two quartic vertices by two internal
lines. Derive the bubble's symmetry factor. Then write its amputated
momentum integral, assigning \(q=p_1+p_2\) through the bubble.
List the other external partitions without evaluating an integral.

\item \textbf{(3 points)} Besides those bubbles, draw the other
connected order-\(\lambda^2\) four-point topology and enumerate its
external assignments. Identify a line whose cut makes it disconnected.
Explain how using full external propagators accounts for this class
and why a logarithm alone does not eliminate it.

\item \textbf{(4 points)} Derive the superficial degree of divergence
\(\omega=dL-2I\), and reduce it in four dimensions to an expression
in \(E\) alone. Apply it to the tadpole and bubble. Explain why a
negative superficial degree for a higher-loop graph does not by
itself guarantee convergence of every subintegration.
\end{parts}

\problem{3}{A mass shift, two regulators, and one physical pole}
\connection{Lecture 22, Sections 1--4 and 7.}
At one loop the scalar self-energy is
\[
 -\ii\Sigma_{\rm tad}=\frac{-\ii\lambda}{2}
 \int\frac{\dd^4k}{(2\pi)^4}\frac{\ii}{k^2-m^2+\ii0}.
\]
You may use the area \(2\pi^2\) of the unit three-sphere and
\(\int\dd^dk_E\,e^{-t k_E^2}/(2\pi)^d=(4\pi t)^{-d/2}\).
\begin{parts}
\item \textbf{(5 points)} Wick rotate the integral, displaying the
factors from the contour and the propagator. Evaluate it with the
Euclidean cutoff \(k_E^2\leq\Lambda^2\), and extract the quadratic
and logarithmic terms at large \(\Lambda\). Check the mass dimension
and state why the result is independent of external momentum.

\item \textbf{(4 points)} Sum repeated 1PI insertions to identify the
pole equation for \(G_2(p)\). Choose the mass parameter to equal the
measured mass \(M\), and determine the one-loop cutoff counterterm
that keeps that pole fixed as \(\Lambda\) changes. Explain what this
says about the regulator dependence of the bare mass.

\item \textbf{(4 points)} Use the Schwinger-parameter representation
of \(1/(k_E^2+m^2)\) to obtain its integral as a gamma function in
\(d\) dimensions. Insert \(\mu^{2\eps}\), set \(d=4-2\eps\), and
expand through the finite term. You may use
\[
 \Gamma(\eps)=\eps^{-1}-\gamma_E+O(\eps),\qquad
 \Gamma(\eps-1)=\frac{\Gamma(\eps)}{\eps-1}.
\]
Explain why the analytically continued answer need not be positive.

\item \textbf{(5 points)} Determine \(\delta m^2\) and
\(\delta Z_\phi\) in \(\MSbar\) at one loop. Derive the relation
between \(M^2\) and \(m^2(\mu)\), including the finite constant.
Use the counterterm vertex to check the sign of the pole cancellation.

\item \textbf{(4 points)} Derive
\(Z_{\rm pole}=[1-\Sigma_R'(M^2)]^{-1}\) and evaluate it at this
order. Compare the finite part of the on-shell mass counterterm with
that of \(\MSbar\). Explain why the two prescriptions do not predict
different measured masses after matching.

\item \textbf{(3 points)} Impose \(\mu\,\dd M^2/\dd\mu=0\) on your
answer in part (d) and determine the leading running of \(m^2(\mu)\).
In taking the derivative, identify which implicit scale dependences
inside the loop term contribute only at higher order.
\end{parts}

\problem{4}{The bubble, running coupling, and two-particle unitarity}
\connection{Lecture 22, Sections 5--8.}
Use on-shell external momenta with pole mass \(M\), so
\(s+t+u=4M^2\). Use \(M\) in the loop integrals as an on-shell mass
organization, and define the coupling in \(\MSbar\). Set
\[
 F(q^2;\mu)=\int_0^1\dd x\log\frac{M^2-x(1-x)q^2-\ii0}{\mu^2}.
\]
No closed form for the real part of \(F\) at general momentum is required.
\begin{parts}
\item \textbf{(4 points)} Prove the Feynman-parameter identity for
\(1/(AB)\). Apply it to the bubble, shift the loop momentum, and use
the Gaussian method of Problem 3 to derive its pole and finite
logarithm. Include the symmetry factor and the phases from both
vertices and propagators.

\item \textbf{(4 points)} Sum the three channels and determine the
\(\MSbar\) coupling counterterm. Write the finite amplitude through
\(O(\lambda^2)\). Explain why its entire pole can be removed by one
momentum-independent \(\phi^4\) counterterm.

\item \textbf{(5 points)} Define \(g\) by
\(\mathcal M(4M^2,0,0)=-g\). Evaluate the parameter integrals at this
point, derive the finite conversion between \(g\) and \(\lambda(\mu)\),
and express the amplitude in terms of \(g,M\) with no arbitrary scale.
Check that the one-loop correction vanishes at the subtraction point.

\item \textbf{(4 points)} Require scale independence of the amplitude
in part (b) to derive \(\beta_\lambda\) at leading order. Solve for
\(\lambda(\mu)\) in terms of \(\lambda(\mu_0)\). Identify the formal
Landau scale and explain the limitation of drawing conclusions there
from a one-loop calculation.

\item \textbf{(4 points)} For \(s>4M^2\), find the interval in \(x\)
where the logarithm has a negative real argument. Determine
\(\operatorname{Im}\mathcal M\), including its sign, using
\(\log(-a-\ii0)=\log a-\ii\pi\). Explain why the \(t\)- and
\(u\)-channel bubbles are real in the physical scattering region.

\item \textbf{(4 points)} Evaluate the integrated two-body phase space
in the center-of-mass frame and verify
\[
 2\operatorname{Im}\mathcal M_{ii}^{(1)}
 =\frac1{2!}\int\dd\Phi_2\,|\mathcal M_{i\to2}^{\rm tree}|^2,
 \quad
 \dd\Phi_2=(2\pi)^4\delta^{(4)}(P-k_1-k_2)
 \prod_{a=1}^2\frac{\dd^3\bm k_a}{(2\pi)^3,2E_a}.
\]
Explain the identical-particle factorial and why a real local
counterterm cannot change this imaginary part.
\end{parts}
\end{document}
